\chapter{ANALISIS DAN INTERPRETASI}
\label{chap:analisis}

Bab ini merangkum makna temuan seluruh \textit{task}. Angka rinci ada pada bab yang dirujuk.

\section{Makna Hasil Pemahaman Data dan Popularitas}
\begin{enumerate}
\item \textit{Venue} adalah seri yang batasnya ditentukan \textit{key}, sehingga \textit{workshop} yang berbagi \textit{key} ikut konferensi induk (7,2\% makalah). Peringkat berlaku untuk definisi ini, dan NeurIPS tetap nomor satu pada definisi lain (Bab~\ref{chap:popularitas}).
\item Data bersih secara struktural, tetapi identitas penulis tidak pasti: nama umum yang menggabungkan banyak orang (\textit{Yang Liu} pada 332 \textit{venue}) menjadi risiko utama pada seluruh hasil \textit{mining}.
\item Aktivitas sangat terkonsentrasi: separuh publikasi berasal dari 5,3\% \textit{venue}, dan 8,3\% penulis menyumbang 52,6\% pasangan makalah-penulis. Porsi sepuluh konferensi teratas naik dari 7,0\% menjadi 11,7\%, karena konferensi kecerdasan buatan dan pemrosesan bahasa alami tumbuh 3,6 sampai 10,2 kali sedangkan konferensi komunikasi dan biomedis menurun. Data tidak memuat penyebabnya.
\end{enumerate}

\section{Makna Pilihan Representasi Transaksi}
\textit{Author-based} menangkap riwayat penulis selama sebelas tahun, sedangkan \textit{author-year-based} hanya menangkap kemunculan pada tahun yang sama. Akibatnya \textit{support} pasangan pada \textit{author-year-based} median 0,29 kali, sekitar 23\% pasangan tidak pernah muncul, dan konferensi dua tahunan berselang seperti ICCV dan ECCV tidak dapat berpasangan. Pilihan representasi karena itu menentukan pola yang dapat ditemukan (Bab~\ref{chap:transaksi}).

\section{Makna Hasil \textit{Mining}}
\begin{enumerate}
\item \textbf{Keterkaitan konferensi mengelompok per bidang.} Kombinasi teratas adalah ICML dan NeurIPS, dan \textit{itemset} terpanjang berisi enam konferensi kecerdasan buatan dan \textit{computer vision} (Bab~\ref{chap:fpm}). Popularitas tidak sama dengan keterkaitan: NeurIPS memiliki \textit{rules} terbanyak, tetapi \textit{lift} \textit{rules} yang menuju kepadanya rendah karena \textit{lift} dibagi \textit{support} konferensi tujuan. Hubungan satu-ke-satu terkuat justru dimiliki pasangan yang lebih kecil dalam satu bidang, misalnya ECOC dan OFC (Bab~\ref{chap:rules}).
\item \textbf{Performa kedua algoritma setara pada data ini.} Transaksi DBLP pendek dan polanya dangkal, sehingga keunggulan FP-Growth yang tidak membangkitkan kandidat belum terlihat pada waktu, tetapi terlihat pada jumlah kandidat Apriori yang tumbuh kuadratik (Bab~\ref{chap:performa}).
\item \textbf{Perhatian bergeser ke kecerdasan buatan.} Kelompok konferensi yang ditemukan tetap sama antarperiode, tetapi pusatnya bergeser dari GLOBECOM dan ICC ke ICLR dan NeurIPS. Sebagian perubahan ditentukan oleh data, yaitu konferensi yang tidak diadakan setiap tahun dan tahun yang tidak tercatat di DBLP (Bab~\ref{chap:temporal}).
\end{enumerate}

\section{Penafsiran yang Perlu Hati-hati}
\begin{enumerate}
\item Nama penulis umum memengaruhi sebagian \textit{itemset} dan \textit{rules} besar, sehingga hasil dibaca bersama keterbatasan pada Bab~\ref{chap:kesimpulan}.
\item \textit{Lift} dipakai untuk membandingkan \textit{rules}, bukan sebagai ukuran mutlak, karena \textit{min\_conf} 50\% dan \textit{support} konferensi yang rendah menjadikan \textit{lift} minimum cukup tinggi.
\item Data 2025 tidak lengkap sehingga periode 2024--2025 pada Bagian E bias.
\end{enumerate}
